% ============================================================
%  第四章 外微分形式与活动标架法
%  结构沿袭苏步青、胡和生《微分几何》第四章（活动标架法，
%  本书的显著特色之一）
%  记号约定：标架向量用 \mathbf{}，微分形式（ω^i、ω_i^j）与
%            各标量系数均不加粗。
% ============================================================

\section{外微分形式}

\subsection{外积}

设 $V$ 是 $\mathbb{R}^3$ 上的向量空间，$(x^1,x^2,x^3)$ 为坐标。微分形式是坐标微元 $\mathrm{d}x^i$ 的反对称多项式。

\begin{definition}[微分形式]
	以下对象分别称为\textbf{$0$ 形式}、\textbf{1 形式}与\textbf{2 形式}：
	\begin{equation}
		0\text{形式：}\omega=f(x^1,x^2,x^3),
		\label{eq:外微分:0形式}
	\end{equation}
	\begin{equation}
		1\text{形式：}\omega=a_1\,\mathrm{d}x^{1}+a_2\,\mathrm{d}x^{2}+a_3\,\mathrm{d}x^{3},
		\label{eq:外微分:1形式}
	\end{equation}
	\begin{equation}
		2\text{形式：}\omega=b_1\,\mathrm{d}x^{2}\wedge\mathrm{d}x^{3}
		+b_2\,\mathrm{d}x^{3}\wedge\mathrm{d}x^{1}
		+b_3\,\mathrm{d}x^{1}\wedge\mathrm{d}x^{2}.
		\label{eq:外微分:2形式}
	\end{equation}
	其中 $\wedge$ 表示\textbf{外积}，满足反对称性：
	\begin{equation}
		\mathrm{d}x^{i}\wedge\mathrm{d}x^{j}=-\mathrm{d}x^{j}\wedge\mathrm{d}x^{i},
		\qquad
		\mathrm{d}x^{i}\wedge\mathrm{d}x^{i}=0.
		\label{eq:外积:反对称}
	\end{equation}
\end{definition}

两个 1 形式 $\omega=\sum_i a_i\mathrm{d}x^{i}$、$\eta=\sum_j b_j\mathrm{d}x^{j}$ 的外积为
\begin{equation}
	\omega\wedge\eta=\sum_{i<j}(a_ib_j-a_jb_i)\,\mathrm{d}x^{i}\wedge\mathrm{d}x^{j}.
	\label{eq:外积:计算公式}
\end{equation}

\begin{remark}
	外积的反对称性式~\eqref{eq:外积:反对称} 正是几何量"定向面积元"与"定向体积元"代数化的来源；两个 1 形式的外积可视为所对应向量场叉积的"对偶"。
\end{remark}

\subsection{外微分}

\begin{definition}[外微分算子]
	外微分 $\mathrm{d}$ 把 $k$ 形式变为 $(k+1)$ 形式，定义如下：
	\begin{equation}
		\mathrm{d}f=\sum_{i}\frac{\partial f}{\partial x^{i}}\,\mathrm{d}x^{i}
		\qquad(0\text{形式}),
		\label{eq:外微分算子:0形式}
	\end{equation}
	\begin{equation}
		\mathrm{d}\omega=\sum_{i}\mathrm{d}a_i\wedge\mathrm{d}x^{i}
		\qquad(\omega\text{为 1 形式}),
		\label{eq:外微分算子:1形式}
	\end{equation}
	并线性地推广到高阶形式。
\end{definition}

\begin{theorem}[外微分的性质]
	外微分满足：
	\begin{enumerate}
		\item 反对称性下的 Leibniz 法则：设 $\omega$ 为 $p$ 形式，则
		\begin{equation}
			\mathrm{d}(\omega\wedge\eta)=\mathrm{d}\omega\wedge\eta+(-1)^{p}\,\omega\wedge\mathrm{d}\eta;
			\label{eq:外微分:Leibniz}
		\end{equation}
		\item $\mathrm{d}^{2}=0$，即任何形式的二阶外微分为零：
		\begin{equation}
			\mathrm{d}(\mathrm{d}\omega)=0;
			\label{eq:外微分:d平方为零}
		\end{equation}
		\item Poincaré 引理：在单连通区域上，闭形式（$\mathrm{d}\omega=0$）必为恰当形式（$\omega=\mathrm{d}\eta$）。
		\label{thm:Poincare引理}
	\end{enumerate}
\end{theorem}

\section{活动标架与结构方程}

\subsection{活动标架}

\begin{definition}[活动标架]
	在 $\mathbb{E}^3$ 中，随点而变的右手正交标架
	\begin{equation}
		\left\{\boldsymbol{r};\;\mathbf{e}_1,\mathbf{e}_2,\mathbf{e}_3\right\}
		\label{eq:活动标架:定义}
	\end{equation}
	称为\textbf{活动标架}。其位移与转动用 1 形式描述：
	\begin{equation}
		\mathrm{d}\boldsymbol{r}=\omega^{1}\mathbf{e}_1+\omega^{2}\mathbf{e}_2+\omega^{3}\mathbf{e}_3,
		\label{eq:活动标架:位移}
	\end{equation}
	\begin{equation}
		\mathrm{d}\mathbf{e}_i=\sum_{j=1}^{3}\omega_{i}^{j}\,\mathbf{e}_j,
		\qquad
		\omega_{i}^{j}=-\omega_{j}^{i}.
		\label{eq:活动标架:转动}
	\end{equation}
	$\omega^{i}$ 称为\textbf{平移形式}，$\omega_{i}^{j}$ 称为\textbf{连接形式}（反对称矩阵 $(\omega_i^j)$）。
\end{definition}

\subsection{结构方程}

\begin{theorem}[Cartan 结构方程]
	活动标架的 1 形式满足下列\textbf{第一结构方程}
	\begin{equation}
		\mathrm{d}\omega^{i}=-\sum_{j=1}^{3}\omega_{j}^{i}\wedge\omega^{j},
		\qquad i=1,2,3,
		\label{eq:结构方程:第一}
	\end{equation}
	与\textbf{第二结构方程}
	\begin{equation}
		\mathrm{d}\omega_{i}^{j}=-\sum_{k=1}^{3}\omega_{k}^{j}\wedge\omega_{i}^{k},
		\qquad i,j=1,2,3.
		\label{eq:结构方程:第二}
	\end{equation}
	它们分别来自二阶微分的可交换性 $\mathrm{d}^{2}\boldsymbol{r}=\boldsymbol{0}$ 与 $\mathrm{d}^{2}\mathbf{e}_i=\boldsymbol{0}$，即 $\mathrm{d}^{2}=0$ 的几何表达。
	\label{thm:Cartan结构方程}
\end{theorem}

\section{活动标架法在曲线论中的应用}

取曲线 $\boldsymbol{r}(s)$（$s$ 为弧长参数）的 Frenet 标架为活动标架：
\begin{equation}
	\mathbf{e}_1=\mathbf{T},\qquad\mathbf{e}_2=\mathbf{N},\qquad\mathbf{e}_3=\mathbf{B}.
	\label{eq:曲线:活动标架}
\end{equation}

\subsection{连接形式与曲率挠率}

由式~\eqref{eq:活动标架:位移} 与 $\left|\boldsymbol{r}'(s)\right|=1$ 得
\begin{equation}
	\omega^{1}=\mathrm{d}s,\qquad\omega^{2}=\omega^{3}=0.
	\label{eq:曲线:平移形式}
\end{equation}

由式~\eqref{eq:活动标架:转动} 与 Frenet 公式（式~\eqref{eq:Frenet公式:方程组}）得
\begin{equation}
	\mathrm{d}\mathbf{e}_1=\mathbf{T}'(s)\,\mathrm{d}s=\kappa\mathbf{N}\,\mathrm{d}s
	\;\Longrightarrow\;
	\omega_{1}^{2}=\kappa\,\mathrm{d}s,\qquad\omega_{1}^{3}=0,
	\label{eq:曲线:连接形式1}
\end{equation}
\begin{equation}
	\mathrm{d}\mathbf{e}_2=\mathbf{N}'(s)\,\mathrm{d}s=\bigl(-\kappa\mathbf{T}+\tau\mathbf{B}\bigr)\mathrm{d}s
	\;\Longrightarrow\;
	\omega_{2}^{1}=-\kappa\,\mathrm{d}s,\qquad\omega_{2}^{3}=\tau\,\mathrm{d}s,
	\label{eq:曲线:连接形式2}
\end{equation}
\begin{equation}
	\mathrm{d}\mathbf{e}_3=\mathbf{B}'(s)\,\mathrm{d}s=-\tau\mathbf{N}\,\mathrm{d}s
	\;\Longrightarrow\;
	\omega_{3}^{2}=-\tau\,\mathrm{d}s,\qquad\omega_{3}^{1}=0.
	\label{eq:曲线:连接形式3}
\end{equation}

因此曲线的曲率 $\kappa$ 与挠率 $\tau$ 恰好出现在连接形式矩阵中：
\begin{equation}
	\bigl(\omega_{i}^{j}\bigr)=
	\begin{pmatrix}
		0 & \kappa & 0\\
		-\kappa & 0 & \tau\\
		0 & -\tau & 0
	\end{pmatrix}\mathrm{d}s.
	\label{eq:曲线:连接形式矩阵}
\end{equation}

\begin{remark}
	这样，曲线论被整理为活动标架的位移 $\mathrm{d}\boldsymbol{r}=\mathbf{e}_1\mathrm{d}s$ 与转动（连接形式）两部分。曲率与挠率刻画了 Frenet 标架沿曲线的转动速率，结构方程自动保证标架保持正交归一。
\end{remark}

\section{活动标架法在曲面论中的应用}

取曲面的切向标架，使 $\mathbf{e}_3=\mathbf{n}$ 为曲面单位法向量，$\mathbf{e}_1,\mathbf{e}_2$ 为切平面上正交单位向量。由于 $\mathrm{d}\boldsymbol{r}$ 恒在切平面上，
\begin{equation}
	\omega^{3}=0,
	\label{eq:曲面:omega3为零}
\end{equation}
于是位移公式式~\eqref{eq:活动标架:位移} 化为
\begin{equation}
	\mathrm{d}\boldsymbol{r}=\omega^{1}\mathbf{e}_1+\omega^{2}\mathbf{e}_2.
	\label{eq:曲面:位移}
\end{equation}

\subsection{第一基本形式与连接形式}

\begin{proposition}[第一基本形式与结构方程]
	第一基本形式为
	\begin{equation}
		\mathrm{I}=\mathrm{d}\boldsymbol{r}\cdot\mathrm{d}\boldsymbol{r}
		=(\omega^{1})^{2}+(\omega^{2})^{2}.
		\label{eq:曲面:第一基本形式}
	\end{equation}
	记 $\omega:=\omega_{1}^{2}$（唯一的独立连接形式），第一结构方程在曲面上化为
	\begin{equation}
		\mathrm{d}\omega^{1}=\omega\wedge\omega^{2},
		\qquad
		\mathrm{d}\omega^{2}=-\omega\wedge\omega^{1}.
		\label{eq:曲面:第一结构方程}
	\end{equation}
	该方程组唯一确定连接形式 $\omega$，它只依赖第一基本形式（即度量），是 Levi--Civita 连接的形式。
\end{proposition}

\subsection{第二基本形式与高斯曲率}

记
\begin{equation}
	a:=\omega_{1}^{3},\qquad b:=\omega_{2}^{3},
	\label{eq:曲面:ab定义}
\end{equation}
它们是刻画曲面弯曲的 1 形式。

\begin{proposition}[第二基本形式]
	由 $\mathrm{d}\mathbf{n}=\mathrm{d}\mathbf{e}_3=\omega_{3}^{1}\mathbf{e}_1+\omega_{3}^{2}\mathbf{e}_2=-a\mathbf{e}_1-b\mathbf{e}_2$ 得第二基本形式
	\begin{equation}
		\mathrm{II}=-\mathrm{d}\boldsymbol{r}\cdot\mathrm{d}\mathbf{n}
		=a\,\omega^{1}+b\,\omega^{2}.
		\label{eq:曲面:第二基本形式}
	\end{equation}
	在切平面上把 $a,b$ 按余标架 $\{\omega^{1},\omega^{2}\}$ 展开：
	\begin{equation}
		a=h_{11}\omega^{1}+h_{12}\omega^{2},\qquad
		b=h_{12}\omega^{1}+h_{22}\omega^{2},
		\label{eq:曲面:ab展开}
	\end{equation}
	则
	\begin{equation}
		\mathrm{II}=h_{11}(\omega^{1})^{2}+2h_{12}\omega^{1}\omega^{2}+h_{22}(\omega^{2})^{2}.
		\label{eq:曲面:第二基本形式展开}
	\end{equation}
\end{proposition}

\begin{theorem}[Gauss 方程与 Codazzi 方程（活动标架形式）]
	由第二结构方程得：
	\textbf{Gauss 方程}
	\begin{equation}
		\mathrm{d}\omega_{1}^{2}=-\omega_{1}^{3}\wedge\omega_{2}^{3}
		=-K\,\omega^{1}\wedge\omega^{2},
		\label{eq:曲面:Gauss方程}
	\end{equation}
	其中高斯曲率为
	\begin{equation}
		K=h_{11}h_{22}-h_{12}^{2};
		\label{eq:曲面:高斯曲率}
	\end{equation}
	\textbf{Codazzi--Mainardi 方程}
	\begin{equation}
		\mathrm{d}\omega_{1}^{3}=-\omega_{2}^{3}\wedge\omega_{1}^{2},\qquad
		\mathrm{d}\omega_{2}^{3}=\omega_{1}^{3}\wedge\omega_{2}^{1}.
		\label{eq:曲面:Codazzi方程}
	\end{equation}
\end{theorem}

\begin{proof}[Gauss 方程的推导]
	由式~\eqref{eq:结构方程:第二} 及 $\omega^{3}=0$、$\omega_{1}^{1}=\omega_{2}^{2}=0$：
	\begin{equation}
		\mathrm{d}\omega_{1}^{2}
		=-\sum_{k}\omega_{k}^{2}\wedge\omega_{1}^{k}
		=-\omega_{3}^{2}\wedge\omega_{1}^{3}.
		\label{eq:曲面:Gauss推导}
	\end{equation}
	又 $\omega_{3}^{2}=-\omega_{2}^{3}=-b$，$\omega_{1}^{3}=a$，故
	\begin{equation}
		\mathrm{d}\omega_{1}^{2}=-(-b)\wedge a=b\wedge a=-a\wedge b.
		\label{eq:曲面:Gauss推导2}
	\end{equation}
	把式~\eqref{eq:曲面:ab展开} 代入并利用 $h_{12}$ 的对称性：
	\begin{equation}
		a\wedge b=(h_{11}h_{22}-h_{12}^{2})\,\omega^{1}\wedge\omega^{2}=K\,\omega^{1}\wedge\omega^{2},
		\label{eq:曲面:Gauss推导3}
	\end{equation}
	即得式~\eqref{eq:曲面:Gauss方程}。
\end{proof}

\begin{theorem}[高斯绝妙定理的又一证明]
	连接形式 $\omega$ 由第一结构方程式~\eqref{eq:曲面:第一结构方程} 唯一确定，而该方程只涉及第一基本形式（度量）。由 Gauss 方程式~\eqref{eq:曲面:Gauss方程}，高斯曲率 $K$ 又由 $\mathrm{d}\omega$ 与面积形式 $\omega^{1}\wedge\omega^{2}$ 之比给出，故 $K$ 是内蕴量，只由第一基本形式决定。
	\label{thm:绝妙定理又一证明}
\end{theorem}

\begin{remark}
	至此，曲面论的三个核心结果——Gauss 方程（含高斯绝妙定理）、Codazzi 方程（曲面论基本定理的可积性条件）——都在活动标架下得到统一的简洁表达。这正是苏步青先生所倡导的 Cartan 活动标架方法的核心思想：用标架的位移与转动来整体刻画几何量。
\end{remark}
